\documentclass[name=Kushaal,title=TwoCanPrelab]{report}
\usepackage{amsmath,graphicx,booktabs,enumitem}
\usepackage[backend=biber,style=numeric]{biblatex}
\addbibresource{ref.bib}
\renewcommand{\name}{Kushaal Singh}
\renewcommand{\reporttitle}{Two-Can Lab Prelab}
\begin{document}
\problem{1: Ideal Flow Coefficients}
From the lecture appendix, with ideal contraction and velocity coefficients $C_c=C_v=1$,
\[
A_c=\frac{\pi D^2}{4},\quad A_o=\frac{\pi d^2}{4},\quad
K_{o,\mathrm{ideal}}=\frac{1}{\sqrt{1-(A_o/A_c)^2}},\quad
K_{\mathrm{ideal}}=K_{o,\mathrm{ideal}}A_o\sqrt{\frac{2g}{A_c}}.
\]
For Can Set 1, use $g=981\,\mathrm{cm/s^2}$ and $1\,\mathrm{in}=2.54\,\mathrm{cm}$.
\begin{center}
\begin{tabular}{lrr}\toprule
Quantity & Top can (1) & Bottom can (2)\\\midrule
$D$ (cm) & 10.17270 & 7.63778\\
$d$ (cm) & 0.63246 & 0.31242\\
$A_c$ (cm$^2$) & 81.2760 & 45.8167\\
$A_o$ (cm$^2$) & 0.314164 & 0.076660\\\bottomrule
\end{tabular}
\end{center}
\[
\boxed{K_1=1.543574\ \mathrm{cm^{3/2}/s}},\qquad
\boxed{K_2=0.501655\ \mathrm{cm^{3/2}/s}}.
\]
These are theoretical values for free discharge; the experimental coefficients can differ.

\problem{2: Sample Simulation}
The supplied fourth-order Runge--Kutta code models
\[
\dot V_1=-K_1\sqrt{\max(V_1,0)},\qquad
\dot V_2=K_1\sqrt{\max(V_1,0)}-K_2\sqrt{\max(V_2,0)}.
\]
For $V_1(0)=406\,\mathrm{mL}$ and $V_2(0)=0$, the predicted bottom-can peak is approximately $255.4\,\mathrm{mL}$ at $19.38\,\mathrm{s}$, corresponding to a depth of $5.57\,\mathrm{cm}$.
\begin{figure}[htbp]
\centering
\includegraphics[width=0.92\linewidth]{two_can_sim_set1_compact.png}
\caption{Simulated water volumes and outlet flow rates for Can Set 1 using ideal coefficients.}
\end{figure}
\clearpage
\problem{3: One-Can Experimental Procedure}
\textbf{Objective.} Determine $K_1$ and $K_2$ separately from emptying experiments, without electronic sensors. For a single can with no inflow,
\[
\dot V=-K\sqrt V,\qquad 2\sqrt{V_0}=KT_e,\qquad K=\frac{2\sqrt{V_0}}{T_e}.
\]
\begin{enumerate}[leftmargin=*,itemsep=0.45em]
\item Set one can upright over a catch reservoir with its bottom orifice temporarily plugged. Keep the can level and ensure the discharge is unobstructed.
\item Measure or confirm its inside diameter $D$ and compute $A_c=\pi D^2/4$. Mark the water heights corresponding to approximately 20\%, 50\%, and 80\% of the can's usable height; calculate $V_0=A_ch_0$ for each mark.
\item Fill to the first marked height. Record the initial water height $h_0$ and volume $V_0$; use the same release technique for every trial.
\item Remove the plug while starting a handheld stopwatch. Stop timing when continuous discharge ends (record the endpoint definition consistently). Record emptying time $T_e$.
\item Repeat each starting height three times, refilling between trials. Record can ID, $D$, $h_0$, $V_0$, $T_e$, and any observations (leaks, splashing, delayed release).
\item Repeat Steps 1--5 for the other can, collecting nine trials per can (18 total).
\item Compute $K=2\sqrt{V_0}/T_e$ for each trial. Plot $2\sqrt{V_0}$ versus $T_e$ and estimate $K$ from the slope of a linear fit. Inspect any nonzero intercept as evidence of timing bias or model mismatch.
\item Report each can's fitted $K$, trial-to-trial variation, and the difference from its theoretical ideal value. Use the experimental coefficients in the two-can simulation for the in-lab prediction.
\end{enumerate}
\textbf{Measurement table columns:} can ID, trial, $D$ (cm), $h_0$ (cm), $V_0$ (mL), $T_e$ (s), $K$ (cm$^{3/2}$/s), notes.

\begin{figure}[htbp]
\centering
\includegraphics[width=0.52\linewidth]{can_dimensions.png}
\caption{Can geometry and Can Set 1 dimensions from the laboratory handout.}
\end{figure}
\end{document}
